\section{Luka analityczna: od stanu stałego do ścieżki przejścia}

Modele równowagi ogólnej są silne tam, gdzie pytanie dotyczy porównania uporządkowanych stanów gospodarki. Z perspektywy zarządu firmy problem BDP jest jednak problemem ścieżki. Firma nie inwestuje w punkcie terminalnym symulacji. Inwestuje w drugim, trzecim i czwartym roku transformacji, gdy nie zna jeszcze finalnego poziomu stóp procentowych, kursu walutowego, popytu, bankructw konkurentów ani własnej zdolności kredytowej.

Raport MFW nie modeluje kilku elementów kluczowych dla decyzji menedżerskiej w Polsce. Po pierwsze, nie opisuje endogenicznej substytucji pracy kapitałem w odpowiedzi na wzrost kosztu pracy i płacy rezerwowej. Po drugie, nie rozróżnia heterogenicznych firm o różnym poziomie marż, dźwigni, płynności i podatności na bankructwo. Po trzecie, nie wprowadza banków jako odrębnych podmiotów z ograniczeniami CAR, LCR, NSFR i portfelem NPL. Po czwarte, w polskim kontekście ważny jest sektor zewnętrzny: kurs PLN/EUR, importochłonność inwestycji, ceny importowane i rachunek bieżący.

Luka nie polega więc na tym, że DSGE jest złe. Luka polega na niedopasowaniu narzędzia do pytania. Jeżeli pytamy o welfare i redystrybucję w stanie równowagi, DSGE jest naturalnym instrumentem. Jeżeli pytamy o trajektorię firm, banków, budżetu państwa i kursu walutowego w trakcie przejścia, potrzebujemy modelu wykonującego przepływy bilansowe w czasie.

Z tego powodu niniejsza praca traktuje BDP jako szok przechodzący przez kilka kanałów transmisji jednocześnie: płacę rezerwową, popyt, inflację, politykę pieniężną, dług publiczny, sektor zewnętrzny, kredyt bankowy i decyzje inwestycyjne firm. W takim ujęciu wynik nie musi być monotoniczny. Ten sam transfer może wzmacniać popyt, ale jednocześnie zwiększać deficyt, koszt kapitału i presję na import.
